% Supplementary material: moved from Sections 2.2, 2.3 and 2.5 of the main text.

\section{Training, Compression, and Deployment Background}
\label{supp:background}

\subsection{Knowledge Distillation and Model Compression}
\label{subsec:background_distillation}
Knowledge distillation, first formalized by \citet{hinton2015distilling}, transfers knowledge from a large ``teacher'' model to a smaller ``student'' by training the student to match the teacher's soft probability distributions (logits) rather than hard labels. This approach preserves the teacher's learned representations of inter-class similarities, information that is lost in one-hot labels, and has proven remarkably effective in the \gls{llm} era. Modern distillation strategies for language models encompass several complementary approaches.

Logit-based distillation trains the student to minimize the KL divergence between its output distribution and the teacher's softened distribution, parameterized by a temperature $\tau$. Google's Gemma 2 \citep{gemma2team2024} and Gemma 3 \citep{team2025gemma3} employ this approach extensively: all Gemma 3 model sizes, including the 27B variant, are trained using knowledge distillation from a larger teacher. The distillation loss is combined with a standard next-token prediction loss, with the relative weighting serving as a key hyperparameter. Feature-based distillation \citep{romero2015fitnets} additionally aligns intermediate representations between teacher and student, capturing structural knowledge that may not be fully expressed in output logits. Microsoft's Orca 2 \citep{mitra2023orca2} demonstrated a variant where a smaller model is taught to use different reasoning strategies depending on task complexity, outperforming models with 5--$10\times$ more parameters on reasoning benchmarks.

Synthetic data generation represents an increasingly dominant form of implicit distillation. Rather than direct logit matching, a powerful \gls{llm} (e.g., GPT-4, Claude 3.5) generates high-quality training examples that are then used to train smaller models. The Phi series pioneered this approach with the ``Textbooks Are All You Need'' methodology \citep{gunasekar2023phi1}, generating synthetic textbook-style content with carefully controlled complexity and diversity. Phi-4's training data is reported to be about 40\% synthetic, generated across 50 dataset types including multi-agent conversations and structured reasoning chains \citep{abdin2024phi4} (\evC); its developers attribute much of its STEM performance to this data, an attribution that has not been tested in a controlled comparison against web-corpus training at equal compute. Post-training quantization, pruning, and adaptation methods are reviewed in \suppsec~\ref{supp:compression}; together with distillation during training they form the usual \gls{slm} optimization pipeline.

\subsection{Efficient Fine-Tuning Methods}
\label{subsec:background_efficient_finetuning}
Full fine-tuning updates all model parameters using task-specific data, which for a 7B-parameter model can require tens of gigabytes of accelerator memory once gradients and optimizer states are included. This cost has made \gls{peft} central to the practical adaptation of \glspl{slm}, particularly when the target deployment environment is a single GPU or an edge-oriented workstation.

\gls{lora} \citep{hu2022lora} addresses this constraint by freezing the pretrained backbone and learning low-rank update matrices for selected linear layers. For a weight matrix $W_0 \in \mathbb{R}^{d \times k}$, the update is parameterized as $\Delta W = BA$, where $B \in \mathbb{R}^{d \times r}$ and $A \in \mathbb{R}^{r \times k}$ with rank $r \ll \min(d,k)$. This drastically reduces the number of trainable parameters while preserving the ability to specialize a model for downstream tasks. \gls{qlora} \citep{dettmers2023qlora} extends this idea by quantizing the frozen base model to 4-bit NormalFloat (NF4) precision and training the adapters in higher precision, allowing models that would ordinarily require multi-GPU setups to be fine-tuned on a single high-memory accelerator.

Beyond LoRA-style adapters, the \gls{peft} landscape includes prompt-based and prefix-based adaptation strategies that shift task information into learned soft prompts or attention prefixes rather than updating the backbone itself. These approaches can be attractive when multiple lightweight task specializations must be maintained simultaneously. However, they tend to underperform LoRA-family methods on complex reasoning tasks, where direct adaptation of attention and feed-forward projections provides a stronger inductive bias for task transfer.

The choice between prompting and fine-tuning therefore depends on task specificity, data availability, and deployment constraints. In-context learning offers zero-training adaptation and is well suited to broad tasks when labeled data are scarce, but it consumes context length and raises inference latency. By contrast, \gls{peft} methods incur an upfront training cost yet reduce per-query overhead and typically yield larger gains on narrow or reasoning-heavy tasks. Prompting and \gls{qlora} results are rarely reported under matched conditions, so they cannot be compared across models from the published record; measuring them is one of the verification studies proposed in Section~\ref{subsec:challenges_verification}.


\subsection{Edge AI and On-Device Inference}
\label{subsec:background_edge_ai}
The edge AI inference stack has matured rapidly, transforming \gls{slm} deployment from a research exercise to a production-ready capability. The llama.cpp framework \citep{gerganov2023llamacpp} is the most widely used runtime for CPU and mixed CPU-GPU inference, and the one against which the others in Table~\ref{tab:deployment_frameworks} position themselves, supporting the \gls{gguf} quantization format with dozens of precision variants (Q2\_K through Q8\_0). Its ecosystem includes llamafile (single-executable packaging), Ollama (user-friendly serving), and extensive platform support spanning x86, ARM, Apple Silicon, and NVIDIA GPUs. Benchmarking on single-board computers shows that llamafile achieves up to $4\times$ higher throughput and 30--40\% lower power consumption than Ollama for equivalent models \citep{nguyen2025sbc}.

Meta's ExecuTorch \citep{meta2024executorch} is an on-device runtime with a small base runtime and backends for CoreML, Qualcomm QNN, ARM, MediaTek, and Vulkan; Meta used it for the on-device measurements of its Llama~3.2 1B and 3B models \citep{meta2024llama32card}. MLC-LLM \citep{mlc2024} deploys through ML compilation (TVM) to iOS, Android, web, and embedded GPUs. ONNX Runtime \citep{microsoft2019onnx} provides cross-platform deployment with optimization passes for diverse hardware. WebLLM \citep{webllm2024} runs inference in the browser with WebGPU and reports retaining up to 80\% of native performance (\evC). Google's LiteRT (formerly TensorFlow Lite) targets mobile CPUs, GPUs, and \glspl{npu} \citep{google2025litert}; its vendor-reported speedups on Qualcomm \glspl{npu} are measured against CPU execution on the vendor's chosen models and are not comparable with the other figures here.

The fundamental constraint in edge deployment is memory bandwidth rather than compute. Flagship phones and mainstream laptops offer on the order of 50--100~GB/s of memory bandwidth, unified-memory laptop designs considerably more, and single-board computers well under 30~GB/s, versus 2--3~TB/s for datacenter GPUs \citep{xu2024device}. Single-stream autoregressive decoding reads the full weight set once per generated token, so in that regime it is bandwidth-bound and quantization from 16 to 4 bits converts a storage saving into a decode-time saving roughly in proportion to the bandwidth freed. The proportionality is regime-dependent rather than general: it weakens under batching as arithmetic intensity rises, largely disappears during the compute-bound prefill phase, and is reduced further on hardware without native low-precision kernels. Section~\ref{subsec:deployment_calc_vs_meas} sets out the conditions in full. Because single-stream interactive use is the dominant on-device regime, INT4 quantization is nevertheless a performance enabler there and not only a storage optimization.


